\honest{Motivated Stopping and Motivated Continuation}{Eliezer Yudkowsky}{2007}

I disagree with some views of the ``Fast and Frugal crowd,'' but they build the most psychologically realistic models of decision. Most experiments hand subjects their options; in real life you must generate them. Likewise you must gather evidence, at a cost, and decide when you have enough. You look at one house and another, adjust how much you need to live near work, and at some point stop and choose. \nb{Search that stops once an option is good enough, with aspirations adjusted along the way, is Herbert Simon's ``satisficing'' (1956).}

Gilovich distinguished motivated skepticism from motivated credulity: for a conclusion we dislike we ask whether the evidence compels it, for one we like whether the evidence allows it. ``I suggest'' an analogous bias in search. With a hidden motive to accept the current best option, we stop and choose too soon; that is motivated stopping. With a hidden motive to reject it, we keep looking and ask for more evidence; that is motivated continuation. \nb{The idea was not new. Yudkowsky had described the stopping half in ``The Third Alternative,'' five months earlier, with the same house example. Kruglanski had argued that motives hasten or delay the end of a search, and Ditto and Lopez, in a 1992 paper titled ``Motivated Skepticism,'' found that people needed less information to reach a conclusion they wanted, and retested an unwelcome medical result more often. The post cites none of this.}

``A major historical scandal in statistics'' was R. A. Fisher's insistence that no causal link between smoking and lung cancer had been established. I say that he testified to Congress: ``Correlation is not causation.'' \nb{Fisher wrote in \textsc{Nature} in 1958 of ``an error \ldots\ of an old kind, in arguing from correlation to causation.'' The editor could find no record that Fisher testified to Congress. An industry consultant presented his hypothesis at the hearings of 1965 and 1969, after Fisher's death.} Or maybe his work as a tobacco consultant gave him a hidden motive to keep looking. He also smoked, and died of colon cancer in 1962. In a footnote I add that Fisher was a frequentist and that Bayesians reason better about causes, citing Judea Pearl. \nb{Pearl holds that causal claims cannot be drawn from associations alone, whether one's probabilities are Bayesian or frequentist.}

Who can argue against gathering more evidence? ``I can.'' Evidence is costly and slow, and refusing to use what you already have is no virtue.

Motivated stopping, I say, appears wherever a third alternative is feared, wherever you would rather not see the obvious counterargument, and wherever you would rather not test a warm glow you paid for. \nb{The post gives no case of stopping. Its one case, Fisher, is of continuation.}

The moral: suspect motivated stopping when you close a search on a comfortable conclusion while fast, cheap evidence remains; suspect motivated continuation when you demand expensive evidence you cannot soon get before doing anything uncomfortable. \nb{Unlike ``The Bottom Line,'' this gives a sign one can check from outside one's own head: the cost of the evidence skipped or demanded.}
